\section{The General Fej\'er--Dirichlet Lift (FD-Lift)}\label{sec:fd-lift-main}

\par\smallskip
\noindent\textit{In this section the FD-lift is formulated, the divisor-sum anchor at integers is established, and the Dirichlet-series factorization $\sum_{n\ge1}\mathcal T_a(n)n^{-s}=\zeta(s)A(s)$ is recorded, together with the polylog–zeta specialization for $\mathfrak S$ and $\mathfrak F$.}
\par\medskip

\subsection*{Standing assumptions and scope}
\noindent Unless stated otherwise, the following assumptions and conventions are used in this section:
\begin{itemize}
  \item If $\sum_{i\ge1}|a(i)|<\infty$, then $\mathcal T_a(z)=\sum_{i\ge1} a(i)\,F(z,i)/i^2$ converges locally uniformly on $\C$ and defines an entire function of order $\le 1$ and exponential type $\le 2\pi$.
  \item If only $\sum_{i\ge1} |a(i)|\,i^{-\sigma}<\infty$ holds for some $\sigma>1$, then Abel-regularized or renormalized versions of the lift are used when entire-ness is required away from the integers (see the renormalized constructions later in the paper).
  \item The Dirichlet-series identity $\sum_{n\ge1}\mathcal T_a(n)n^{-s}=\zeta(s)A(s)$ is asserted on the half-plane of absolute convergence of $A(s)=\sum_{n\ge1}a(n)n^{-s}$ (and extended by analytic continuation when available).
\end{itemize}

\begin{quote}\small
\textbf{Bridge at a glance.} For weights $a:\mathbb N\to\mathbb C$,
\[
\text{local:}\quad \mathcal T_a(n)=(a*1)(n)
\quad \text{corresponds to} \quad
\text{spectral:}\quad \sum_{n\ge1}\frac{\mathcal T_a(n)}{n^s}=\zeta(s)\,A(s),
\]
where $A(s)=\sum_{n\ge1}a(n)n^{-s}$. The kernel $F(z,i)/i^2$ acts as a divisor filter at integers, while Dirichlet convolution transfers to multiplication in the Dirichlet–series domain.
\end{quote}

This section introduces the general framework of the Fej\'er--Dirichlet Lift (FD-Lift). The method constructs an entire function $\mathcal{T}_a(z)$ from a given arithmetic sequence of weights $a(n)$. This function serves as an analytic interpolant for the Dirichlet convolution $(a*1)(n)$ and possesses a Dirichlet series that factors in terms of the Riemann Zeta function.
This section abstracts the mechanism of Section~\ref{sec:application-prime-indicator} into a linear operator $\mathcal{T}_a$.

\subsection{Definition and Core Properties}

The construction of the lift is a linear superposition of the fundamental Fej\'er kernel terms, weighted by the arithmetic information encoded in a sequence $a:\N \to \C$. The logic follows from the observation that the term $F(z,i)/i^2$ serves as an analytic continuation of the divisor-indicator function $\mathbf{1}_{i|n}$.

\begin{definition}[The General Fej\'er--Dirichlet Lift]
For a complex sequence of weights $a:\N \to \C$, the FD-Lift $\mathcal{T}_a(z)$ is defined as
\[
\mathcal{T}_a(z) := \sum_{i=1}^{\infty} \frac{a(i)}{i^2} F(z,i),\qquad F(z,1):=1.
\]
\end{definition}

If the series of weights converges absolutely, $\sum_{i=1}^{\infty} |a(i)| < \infty$, the series defining $\mathcal{T}_a(z)$ converges uniformly on compact subsets of $\C$. Since each term in the sum is an entire function, the limit function $\mathcal{T}_a(z)$ is, by the Weierstrass theorem on uniformly convergent series of analytic functions, an entire function of order at most $1$ and exponential type at most $2\pi$.

The constructed function $\mathcal{T}_a(z)$ possesses two foundational properties that establish its role as a bridge between the local space of integers and the spectral space of Dirichlet series. These are summarized in the following theorem.

\begin{lemma}[Growth bound and exponential type]\label{lem:growth-type}
For $z=x+\mathrm{i}y\in\C$ and integer $i\ge2$ the Fej\'er–Dirichlet kernel satisfies
\[
\bigl|F(z,i)\bigr|\;\le\; i^2\,\cosh(2\pi|y|).
\]
Consequently, if $\sum_{i\ge1}|a(i)|<\infty$, then the FD-Lift
\[
\mathcal{T}_a(z)\;=\;\sum_{i\ge1} a(i)\,\frac{F(z,i)}{i^2}
\]
obeys the global bound
\[
\bigl|\mathcal{T}_a(x+\mathrm{i}y)\bigr|\;\le\;\Bigl(\sum_{i\ge1}|a(i)|\Bigr)\,\cosh(2\pi|y|),
\]
and in particular $\mathcal{T}_a$ is an entire function of order at most $1$ and exponential type at most $2\pi$, since $\cosh(2\pi|y|)\le \tfrac12(e^{2\pi|y|}+e^{-2\pi|y|})$ provides a global $e^{2\pi|y|}$–bound along vertical lines.
\end{lemma}

\begin{remark}[Order $\le 1$ and exponential type $\le 2\pi$]
An entire function $H$ has exponential type at most $T$ if for every $\epsilon>0$ there exists $C_\epsilon$ with $|H(z)|\le C_\epsilon\,e^{(T+\epsilon)|z|}$ for all $z\in\C$ (see, e.g., Levin~\cite{levin1996}). A bound $|H(x+\ii y)|\le C\,e^{T|y|}$ holding uniformly for all $x,y\in\mathbb R$ is sufficient for this, since $|y|\le|z|$. Lemma~\ref{lem:growth-type} yields 
\[
|\mathcal T_a(x+\ii y)|\le \Bigl(\sum |a(i)|\Bigr)\cosh(2\pi|y|)\le C\,e^{2\pi|y|}
\qquad(x,y\in\mathbb R).
\]
Since $|y|\le |z|$ for $z=x+\ii y$, this immediately gives the global bound
\[
|\mathcal T_a(z)|\le C\,e^{2\pi|z|}\qquad(z\in\mathbb C),
\]
so $\mathcal T_a$ is of order at most $1$ and exponential type at most $2\pi$ in the usual radial sense. The geometric corrector $z\mapsto q^{-z}=e^{-z\log q}$ (used elsewhere in the paper) is not covered by this argument: it satisfies $|q^{-z}|=e^{-\Re(z\log q)}\le e^{|\log q|\,|z|}$, has exponential type $|\log q|$, and is unbounded on $\mathbb R$. Functions containing it are treated separately in Theorem~\ref{thm:analyticity-growth} and Lemma~\ref{lem:fsharp-growth-corrected}.
\end{remark}

\begin{proof}
Using the cosine series $F(z,i)= i+2\sum_{k=1}^{i-1}(i-k)\cos\!\bigl(2\pi k z/i\bigr)$ and $|\cos(x+\mathrm{i}y)|\le\cosh(y)$ gives
\[
|F(z,i)|\;\le\; i\cdot\cosh(0)\;+\;2\sum_{k=1}^{i-1}(i-k)\cosh\!\Bigl(\tfrac{2\pi k}{i}|y|\Bigr)
\;\le\; i\;+\;2\sum_{k=1}^{i-1}(i-k)\cosh(2\pi|y|).
\]
Since $\sum_{k=1}^{i-1}(i-k)=\tfrac{i(i-1)}{2}$, it follows that $|F(z,i)|\le i^2\cosh(2\pi|y|)$. Summation with weights $|a(i)|/i^2$ yields the claimed bound for $\mathcal{T}_a$. The growth estimate implies order $\le 1$ and exponential type $\le 2\pi$.
\end{proof}

\begin{theorem}[Fundamental Properties of the FD-Lift]\label{thm:fd-lift-core}
Let $\mathcal{T}_a(z)$ be the FD-Lift for a sequence of weights $a(n)$ whose associated Dirichlet series is $A(s) = \sum_{n=1}^\infty a(n)n^{-s}$. The following properties hold:
\begin{enumerate}
    \item \textbf{(Local Property: Integer Values)} For any integer $n\ge1$, the value of the function is equal to the Dirichlet convolution of $a$ and the unit function $1(n)=1$ (at integer arguments the defining series reduces to the finite sum over the divisors of $n$, so no summability assumption is needed for this part):
    \[
    \mathcal{T}_a(n) = \sum_{d|n} a(d) = (a*1)(n).
    \]
    \item \textbf{(Spectral Property: Dirichlet Series)} In the half-plane of absolute convergence, the Dirichlet series of the sequence $\mathcal{T}_a(n)$ is the product of the Riemann Zeta function and the Dirichlet series of the weights:
    \[
    \sum_{n=1}^{\infty}\frac{\mathcal{T}_a(n)}{n^s} = \zeta(s)A(s).
    \]
    This identity holds for all $s$ with $\Re s>\max\{1,\sigma_0\}$, where $A(s)=\sum_{n\ge1} a(n)n^{-s}$ converges absolutely on $\Re s>\sigma_0$; extension to other $s$ follows by analytic continuation, when available.
\end{enumerate}
\end{theorem}
\begin{proof}[Sketch]
(1) For integers $n$, the projector property yields $F(n,i)/i^2=\mathbf{1}_{i\mid n}$, hence $\mathcal T_a(n)=(a*1)(n)$.
(2) For the Dirichlet series, swap sums (absolute convergence) and regroup: $\sum_n \mathcal T_a(n)n^{-s}
= \sum_d a(d)d^{-s}\sum_m m^{-s} = A(s)\zeta(s)$.
Details: Appendix~\ref{app:fdlift-proofs}.
\end{proof}

\begin{figure}[!htbp]
\centering
\includegraphics[width=\textwidth]{plot_mu_interpolant.pdf}
\caption{The FD-Lift for the M\"obius function, $\mathcal{T}_\mu(x)$. The continuous function (blue curve), generated from the sum of Fej\'er kernels weighted by $\mu(i)$, is shown to correctly interpolate the discrete values of the convolution $(\mu*1)(n)$ (red dots). This provides a visual confirmation of Theorem~\ref{thm:fd-lift-core} and illustrates the construction of an "analytic annihilator" for integers $n>1$.}
\label{fig:mu_interpolant}
\end{figure}

\FloatBarrier 

\begin{corollary}[Order $1$ and exponential type $\le 2\pi$]
Assume $\sum_{i=1}^{\infty} |a(i)| < \infty$. Then the FD-Lift $\mathcal{T}_a$ is an entire function of order at most $1$ and exponential type at most $2\pi$. In particular, there exists a constant $C=C(a)>0$ such that
\[
|\mathcal{T}_a(x+\ii y)| \;\le\; C\, e^{2\pi |y|} \qquad (x,y\in\R),
\]
and hence, since $|y|\le |x+\ii y|=|z|$, also the radial bound
\[
|\mathcal{T}_a(z)| \;\le\; C\, e^{2\pi |z|} \qquad (z\in\C)
\]
holds. Consequently, the order is $\le 1$ and the exponential type is $\le 2\pi$.
\end{corollary}


\subsection{Functional and Algebraic Properties}

The FD-Lift framework respects the fundamental algebraic structures of arithmetic functions and their associated Dirichlet series. The lift operator $\mathcal{T}: a \mapsto \mathcal{T}_a$ exhibits the following key properties.

\begin{itemize}
    \item \textbf{Linearity.} The lift is a linear operator. For any two sequences $a_1, a_2$ and complex constants $c_1, c_2$, it holds that
    \[ \mathcal{T}_{c_1 a_1 + c_2 a_2}(z) = c_1 \mathcal{T}_{a_1}(z) + c_2 \mathcal{T}_{a_2}(z). \]

    \item \textbf{M\"obius Inversion.} Since $\mathcal{T}_a(n) = (a*1)(n)$, the original weight sequence $a(n)$ can be recovered from the integer values of the lift by Dirichlet convolution with the M\"obius function $\mu(n)$:
    \[ a(n) = (\mu * \mathcal{T}_a)(n). \]
    In the spectral space, this corresponds to the division of the resulting Dirichlet series by $\zeta(s)$.

    \item \textbf{Logarithmic Derivative Identity.} A relationship holds between the logarithmic derivatives of the involved Dirichlet series. Let $M(s)=\zeta(s)A(s)$ denote the spectrum of the lift. Then
    \[
    -\frac{M'(s)}{M(s)} \;=\; -\frac{\zeta'(s)}{\zeta(s)} \;-\; \frac{A'(s)}{A(s)}.
    \]
    This is a purely analytic identity. An arithmetical “von Mangoldt” interpretation of $-A'(s)/A(s)$ requires that $A(s)$ admit an Euler product (e.g., when the weights are multiplicative and satisfy standard convergence hypotheses); without such a product, $-A'(s)/A(s)$ is only a formal logarithmic derivative and does not carry canonical positivity or prime-power counting properties.

    \item \textbf{Divisor lattice algebra.} The basis functions $F(z,i)/i^2$ encode divisibility at integer arguments. For integers $n,i,j$ one has the pointwise identity
    \[
    \frac{F(n,i)}{i^2}\,\frac{F(n,j)}{j^2}=\mathbf{1}_{i\mid n}\,\mathbf{1}_{j\mid n}=\mathbf{1}_{\mathrm{lcm}(i,j)\mid n}=\frac{F(n,\mathrm{lcm}(i,j))}{\mathrm{lcm}(i,j)^2},
    \]
    which highlights compatibility with Dirichlet convolution. No such identity is claimed for non-integer arguments.

    \item \textbf{Stationarity at the integers.} If $\sum_{i\ge1}|a(i)|<\infty$, then $\mathcal T_a'(n)=0$ for every integer $n\ge1$: each $\phi_i$ is stationary at every integer (Lemma~\ref{lem:int-values}), and the differentiated series converges uniformly because $|\phi_i'|\le2\pi$ (Lemma~\ref{lem:global-derivative-bound}). By Lemma~\ref{lem:second-derivatives} the same termwise argument gives, for every integer $n\ge1$,
    \[
    \mathcal T_a''(n)\;=\;-\frac{2\pi^2}{3}\sum_{i\mid n}a(i)\Bigl(1-\frac1{i^2}\Bigr)\;+\;2\pi^2\sum_{i\nmid n}\frac{a(i)}{i^2\sin^2(\pi n/i)} ,
    \]
    and the series converges absolutely, since the $i$-th summand of the second sum is bounded by $2\pi^2|a(i)|/(i^2\sin^2(\pi n/i))\le\pi^2|a(i)|/2$ for $i\nmid n$ (Jordan's inequality, as in the proof of Lemma~\ref{lem:K-bounds}). In words: an FD-lift with absolutely summable weights passes through every integer value with zero slope. This is the origin of the \emph{tangent mismatch} between the Fej\'er sum $S_q$ and the exponential corrector discussed in Section~\ref{sec:periodic-normalizer}.
\end{itemize}

\subsection{Connection to Prime Number Theory: Two Complementary Approaches}

A key application of the FD-Lift is its ability to establish a direct connection to prime number theory and the zeros of the Riemann Zeta function. Two fundamentally different but complementary approaches to leverage this connection are presented.
\begin{enumerate}
    \item \textbf{The Direct Adapter Method (Analysis).} This is a targeted approach where the weights $a(n)$ are precisely tailored to isolate a specific arithmetic property---primality, as represented by the von Mangoldt function $\Lambda(n)$---and thereby directly generate its spectrum, $-\zeta'(s)/\zeta(s)$.
    \item \textbf{The Dynamic Method (Synthesis).} This is a synthetic approach where, instead of isolating one property, a general, parameter-dependent family of functions (related to $\sigma_{-s}(n)$) is constructed. A new, complex property related to prime numbers is then derived through an analytic operation (differentiation) within this family.
\end{enumerate}
Both tracks illuminate the same underlying structures from different perspectives and demonstrate the generality and flexibility of the FD-Lift framework. The following sections detail each approach.

\begin{quote}\small
\textbf{Takeaways (FD-lift).}
Integer anchors satisfy $\mathcal T_a(n)=\sum_{d\mid n}a(d)$.
The Dirichlet-series factorization $\sum_{n\ge1}\mathcal T_a(n)n^{-s}=\zeta(s)A(s)$ is obtained.
For $a(i)=(q-1)q\,q^{-i}$ the specializations for $\mathfrak S$ and $\mathfrak F$ follow via $\operatorname{Li}_s(1/q)$.
\end{quote}


\subsection{Track 1: The Direct Adapter Method for Isolating \texorpdfstring{$\Lambda(n)$}{Lambda(n)}}

\par\smallskip
\noindent\textit{In this subsection a direct adapter is set up to isolate the von Mangoldt function at integers via the FD-lift. A renormalized kernel is employed to ensure absolute and locally uniform convergence; the integer anchor $\mathcal T^{(\mathrm{ren})}_{\mu*\Lambda}(n)=\Lambda(n)$ is verified and the associated Dirichlet-series mechanism is recorded.}
\par\medskip

\begin{quote}\small
\textbf{Example (three lines).}
For $a=\mu*\Lambda$, a naive lift $\sum_{i\ge1}(\mu*\Lambda)(i)\,F(z,i)/i^2$ fails to be absolutely convergent and is unstable under Abel summation.
A renormalized difference kernel $(\phi_i-\phi_\infty)$ is therefore used; this guarantees absolute and locally uniform convergence while preserving $\mathcal T^{(\mathrm{ren})}_{\mu*\Lambda}(n)=\Lambda(n)$.
The construction is detailed in Appendix~\ref{app:fdlift-proofs}, Theorem~\ref{thm:renormalized-muLambda}.
\end{quote}

\begin{proposition}[Explicit formula for $\psi$ (classical), in the normalization used here]\label{prop:adapter-explicit}
Let $x>1$ and let $\psi(x)=\sum_{n\le x}\Lambda(n)$ be Chebyshev's function in the symmetric sense. Then
\[
\psi(x)= x - \sum_{\rho}\frac{x^{\rho}}{\rho} - \log(2\pi) - \frac{1}{2}\log\!\bigl(1-x^{-2}\bigr),
\]
where the sum runs over the nontrivial zeros $\rho$ of $\zeta(s)$ with symmetric summation.
\end{proposition}

\begin{remark}[Explicit formula: form used]
Throughout, the explicit formula is understood in the form of Proposition~\ref{prop:adapter-explicit}, including the contribution of the trivial zeros via $-\tfrac12\log(1-x^{-2})$ and the constant $-\log(2\pi)$. All later references point to this exact normalization.
\end{remark}

\begin{proof}
Fix $c>1$ and $T\ge2$. Consider
\[
I_T:=\frac{1}{2\pi i}\int_{c-iT}^{c+iT} -\frac{\zeta'(s)}{\zeta(s)}\,\frac{x^{s}}{s}\,ds.
\]
By Perron's formula in its symmetric form (see, e.g., \cite[Thms.\ 4.11 and 4.16]{titchmarsh1986}; see also \cite[Ch.\ 3]{edwards2001} and \cite{iwanieckowalski2021}), $I_T$ tends to $\psi(x)$ as $T\to\infty$, with the standard half-weight convention at integers. For $x>1$ and $x\notin\mathbb N$ the same limit holds without half-weights. Shift the contour to the left across $s=1$, the nontrivial zeros $\rho$, $s=0$, and the trivial zeros $s=-2k$ ($k\in\mathbb{N}$). The residues are:
\begin{itemize}
\item at $s=1$: $x$;
\item at $s=\rho$: $-\dfrac{x^{\rho}}{\rho}$;
\item at $s=0$: $-\dfrac{\zeta'(0)}{\zeta(0)}=-\log(2\pi)$;
\item at $s=-2k$: $-\dfrac{x^{-2k}}{-2k}=\dfrac{1}{2}\dfrac{x^{-2k}}{k}$, summing to $-\dfrac{1}{2}\log(1-x^{-2})$.
\end{itemize}
Along a sequence $T\to\infty$ chosen so that $T$ stays away from the ordinates of the nontrivial zeros (the usual choice, keeping the distance $\gg1/\log T$), the horizontal segments and the left vertical segment vanish by the standard bounds for $\zeta$ and $\zeta'/\zeta$ on such lines; see \cite{edwards2001,titchmarsh1986,iwanieckowalski2021}. The argument is classical and does not use the adapter; the adapter enters only through the identification of $-\zeta'/\zeta$ as the spectrum of $\mathcal T^{(\mathrm{ren})}_{\mu*\Lambda}$. Summing residues yields the stated identity.
\end{proof}

This approach provides the most direct path to embedding the pole structure of the Zeta function's logarithmic derivative into the spectrum of an entire function.

\subsubsection{Strategy: From Zeros to Poles}
A central goal in analytic number theory is to understand the distribution of the non-trivial zeros $\rho$ of the Zeta function. In complex analysis, poles are often more tractable than zeros. The strategy is therefore to construct a function whose poles coincide exactly with the zeros of $\zeta(s)$. The canonical function with this property is the logarithmic derivative, $-\zeta'(s)/\zeta(s)$. On the $z$-side, the renormalized entire interpolant from Theorem~\ref{thm:renormalized-muLambda} is used, thereby avoiding unrenormalized Abel-limit issues away from integers.
The FD-Lift is "tuned" by choosing weights $a(n)$ such that the resulting spectrum matches this target:
\[
\sum_{n=1}^{\infty}\frac{\mathcal T_a(n)}{n^s} = \zeta(s) A(s) \stackrel{!}{=} -\frac{\zeta'(s)}{\zeta(s)}.
\]

\subsubsection{The Adapter Mechanism}
To achieve the target spectrum, the Dirichlet series of the weights must be $A(s) = -\frac{\zeta'(s)}{\zeta(s)^2}$. From the theory of Dirichlet series, the arithmetic function corresponding to this spectrum is the Dirichlet convolution of the M\"obius function $\mu$ and the von Mangoldt function $\Lambda$. The mechanism proceeds in three steps:
\begin{enumerate}
    \item \textbf{Define the weights.} The weights are set to $a(n) := b(n) = (\mu * \Lambda)(n)$. The associated Dirichlet series is then, as required:
    \[
    B(s) = \sum_{n=1}^{\infty}\frac{b(n)}{n^s} = \left(\sum\frac{\mu(n)}{n^s}\right)\left(\sum\frac{\Lambda(n)}{n^s}\right) = \frac{1}{\zeta(s)}\cdot\left(-\frac{\zeta'(s)}{\zeta(s)}\right) = -\frac{\zeta'(s)}{\zeta(s)^2}.
    \]
    \item \textbf{Evaluate the lift (spectral side).} The spectrum of the lift $\mathcal{T}_b(z)$ is, by Theorem~\ref{thm:fd-lift-core},
    \[
    \sum_{n=1}^{\infty}\frac{\mathcal T_b(n)}{n^s} = \zeta(s)B(s) = \zeta(s) \left(-\frac{\zeta'(s)}{\zeta(s)^2}\right) = -\frac{\zeta'(s)}{\zeta(s)}.
    \]
    \item \textbf{Evaluate the lift (local side).} The value at integer arguments is $\mathcal{T}_b(n)=(b*1)(n) = (\mu * \Lambda * 1)(n)$. Using the standard identities $(\Lambda*1)(n)=\log n$ and $(\mu*\log)(n)=\Lambda(n)$ (see, e.g., Apostol \cite[Ch.~2]{apostol1976}), it follows that:
    \[
    \mathcal{T}_{(\mu*\Lambda)}(n) = \Lambda(n).
    \]
\end{enumerate}

\begin{theorem}[Renormalized FD-lift for $a=\mu*\Lambda$]\label{thm:renormalized-muLambda}
Let $\phi_i(z):=F(z,i)/i^2$ and $\phi_\infty(z):=(\sin(\pi z)/(\pi z))^2$ with $\phi_\infty(0):=1$. Then the renormalized series
\[
\mathcal{T}^{(\mathrm{ren})}_{\mu*\Lambda}(z)\ :=\ \sum_{i=1}^{\infty}(\mu*\Lambda)(i)\,\bigl(\phi_i(z)-\phi_\infty(z)\bigr)
\]
converges absolutely and locally uniformly on $\C$ and hence defines an entire function. Moreover, for every integer $n\ge1$,
\[
\mathcal{T}^{(\mathrm{ren})}_{\mu*\Lambda}(n)=\Lambda(n).
\]
\end{theorem}
\begin{proof}[Sketch]
On $|z|\le R$: $\sin(\pi z/i)=\frac{\pi z}{i}(1+E_i(z))$, hence $F(z,i)/i^2=\phi_\infty(z)(1+O_R(i^{-2}))$.
With $(\mu*\Lambda)(i)\ll\log i$ this gives absolute local uniform convergence.
At integers, the renormalization removes the $i=\infty$ limit and yields $\Lambda(n)$.
Full proof: Appendix~\ref{app:fdlift-proofs}.
\end{proof}


\begin{remark}[Abel regularisation away from integers]
For $0<r<1$ the unrenormalized Abel sum $\sum_{i\ge1}(\mu*\Lambda)(i)\,\phi_i(z)\,r^i$ splits as $\sum (\mu*\Lambda)(i)\,(\phi_i-\phi_\infty)(z)\,r^i+\phi_\infty(z)\sum (\mu*\Lambda)(i)\,r^i$. The first term converges locally uniformly as $r\uparrow1$, but the scalar series $\sum (\mu*\Lambda)(i)\,r^i$ need not have a finite Abel limit due to the pole structure of $(1/\zeta)'(s)$ at nontrivial zeros of $\zeta$. Hence existence of the unrenormalized Abel limit for $z\notin\mathbb Z$ is not asserted.
\end{remark}

\subsubsection{Resulting Pole Structure and Interpretation}
For $0<r<1$ define
\[
\mathcal{T}^{(\mathrm{Ab})}_{(\mu*\Lambda)}(z;r):=\sum_{i=1}^{\infty} (\mu*\Lambda)(i)\,\frac{F(z,i)}{i^2}\, r^i,
\]
which is entire in $z$ for each fixed $r$. At integer arguments $n$ the Abel sum stabilizes to $\Lambda(n)$ as $r\uparrow1$ because only finitely many terms contribute. For $z\notin\mathbb Z$ existence of the unrenormalized Abel limit is not asserted; instead the renormalized entire interpolant from Theorem~\ref{thm:renormalized-muLambda} is used. The associated Dirichlet series $-\zeta'(s)/\zeta(s)$ has simple poles at:
\begin{itemize}
    \item $s=1$, with residue $+1$, corresponding to the main term $x$ in the Prime Number Theorem,
    \item all nontrivial zeros $\rho$ of $\zeta(s)$, with residue $-m_\rho$ (where $m_\rho$ is the multiplicity of the zero),
    \item the trivial zeros $s=-2k$, with residue $-1$.
\end{itemize}

\subsubsection{Explicit formula via Perron}
The explicit formula for $\psi(x)=\sum_{n\le x}\Lambda(n)$ stated in Proposition~\ref{prop:adapter-explicit} applies here verbatim and is recalled for convenience:
\[
\psi(x)\;=\; x \;-\; \sum_{\rho}\frac{x^{\rho}}{\rho} \;-\; \log(2\pi) \;-\; \tfrac{1}{2}\log\!\Bigl(1-\frac{1}{x^{2}}\Bigr).
\]

\begin{corollary}[RH]
Under the Riemann Hypothesis one has
\[
\psi(x) = x + O\!\big(x^{1/2}\log^{2}x\big).
\]
This is classical and goes back to von Koch~\cite{vonkoch1901}; it is quoted here only for orientation and is not used elsewhere in this paper. Textbook proofs are in \cite{edwards2001,titchmarsh1986}.
\end{corollary}

The FD-Lift thus provides a constructive perspective on the appearance of the Zeta zeros in explicit formulas for primes: they arise as the poles of the spectrum of the entire function that naturally interpolates the prime-power detecting function $\Lambda(n)$. The result of this mechanism is visualized in Figure~\ref{fig:lambda_interpolant}.

\begin{figure}[!htbp]
\centering
\includegraphics[width=\textwidth]{plot_lambda_interpolant.pdf}
\caption{Visual confirmation of the direct adapter method. The smooth, entire function $\mathcal{T}_{(\mu*\Lambda)}(x)$ (red curve) is shown passing exactly through the discrete values of the von Mangoldt function $\Lambda(n)$ (grey stems). This illustrates how the FD-Lift provides a continuous interpolant for the prime-power detecting function.}
\label{fig:lambda_interpolant}
\end{figure}

\FloatBarrier

\begin{quote}\small
\textbf{Takeaways (Track 1: direct adapter for $\Lambda$).}
A renormalized adapter based on $(\phi_i-\phi_\infty)$ is used to enforce absolute/local uniform convergence.
At integers the interpolation equals $\Lambda(n)$.
Its spectrum is the classical Dirichlet series $-\zeta'(s)/\zeta(s)=\sum_{n\ge1}\Lambda(n)n^{-s}$.
\end{quote}

\subsection{Track 2: The Dynamic Method via the Two-Variable Lift}\label{sec:dynamic-method}

\par\smallskip
\noindent\textit{In this subsection a two-variable FD-lift $\mathcal T^{(\mathrm{ren})}(z,s)$ is introduced; holomorphy on $\Re s>-1$ is documented, the integer identity $\mathcal T^{(\mathrm{ren})}(n,s)=\sum_{d\mid n}d^{-s}$ is established, and differentiation at $s=0$ yields $-\sum_{d\mid n}\log d=-(\tau*\Lambda)(n)$, whose Dirichlet series is $\zeta(u)\zeta'(u)$.}
\par\medskip


\noindent\textbf{Variables.} Throughout this section, $z$ denotes the interpolation variable, $s$ the weight parameter in the two--variable lift, and $u$ the Dirichlet--series variable. This convention prevents clashes between the interpolation and spectral sides.

This second approach explores a dynamic structure by introducing a new complex parameter into the weights themselves. Instead of constructing a single function, an entire family of functions is generated, and a new arithmetic relationship is revealed by analyzing the behavior of this family.

\subsubsection{A Parameter-Dependent Family of Functions}
The lift is generalized by choosing weights that are themselves a function of a complex variable $s$, namely $a(i,s) = i^{-s}$. This elevates the FD-Lift to a function of two complex variables, $\mathcal{T}(z,s)$.

\begin{definition}[The Two-Variable Lift]
For $z \in \C$ and $s \in \C$, the two-variable lift is defined as
\[\mathcal{T}(z,s) := \sum_{i=1}^{\infty} \frac{i^{-s}}{i^2} F(z,i) = \sum_{i=1}^{\infty} \frac{F(z,i)}{i^{s+2}}.\]
\end{definition}

\noindent\textbf{Convergence and regularity.}
For each fixed $s$ with $\Re s>1$, the series defining $\mathcal{T}(z,s)$ converges absolutely and uniformly on compact subsets of $z\in\C$, hence $z\mapsto\mathcal{T}(z,s)$ is entire.
For each fixed $z\in\C$, the series converges absolutely for $\Re s>1$ and defines a function holomorphic in $s$ on $\Re s>1$.
Moreover, gliding differentiation with respect to $s$ is justified on every half-plane $\Re s\ge 1+\delta$ with $\delta>0$ by the Weierstrass $M$-test applied to the majorant $\sum_{i\ge1}(\log i)/i^{1+\delta}$.
At integer points $n\in\mathbb{N}$ one has the interpolation identity
\[
\mathcal{T}(n,s)=\sum_{d\mid n} d^{-s}=\sigma_{-s}(n).
\]

\begin{itemize}
    \item \textbf{In the local space}, at an integer argument $n$, each function in the family interpolates the sum-of-powers-of-divisors function $\sigma_{-s}(n)$:
    \[
    \mathcal{T}(n,s) = \sum_{d|n} d^{-s} = \sigma_{-s}(n).
    \]
    \item \textbf{In the spectral space}, the Dirichlet series of this sequence (over the variable $u$) is the product of two Zeta functions, a central object in number theory:
    \[
    \sum_{n=1}^{\infty}\frac{\sigma_{-s}(n)}{n^u} = \zeta(u)\zeta(u+s), \quad \text{for } \re(u)>1, \re(u+s)>1.
    \]
\end{itemize}
The parameter $s$ acts as a "dial" that continuously adjusts the arithmetic properties of the system.

\subsubsection{The Differentiation Operator as an Arithmetic Transformer}
The potential of this dynamic approach is demonstrated by applying an analytic operation to the parameter $s$: differentiation at the origin $s=0$. The \emph{differentiation operator} is defined as $\mathcal{D} := \frac{\partial}{\partial s}\big|_{s=0}$. It transforms the arithmetic properties of the function on all levels of the lift.

\begin{proposition}[Renormalized spectral differentiation at $s=0$]\label{prop:spectral-diff}
Let $\phi_i(z):=F(z,i)/i^2$ and $\phi_\infty(z):=\bigl(\frac{\sin(\pi z)}{\pi z}\bigr)^2$ with $\phi_\infty(0):=1$.
For $\Re s>-1$ define the renormalized two-variable lift
\[
\mathcal T^{\mathrm{ren}}(z,s)\ :=\ \sum_{i\ge1}\bigl(\phi_i(z)-\phi_\infty(z)\bigr)\,i^{-s}.
\]
For $\Re s>1$ it is related to the two-variable lift by $\mathcal T(z,s)=\mathcal T^{\mathrm{ren}}(z,s)+\phi_\infty(z)\,\zeta(s)$; this provides the meromorphic continuation of $s\mapsto\mathcal T(z,s)$ to $\Re s>-1$, with at most a simple pole at $s=1$ of residue $\phi_\infty(z)$.
Then:
\begin{enumerate}
\item For each fixed $s$ with $\Re s>-1$, the map $z\mapsto \mathcal T^{\mathrm{ren}}(z,s)$ is entire; for each fixed $z\in\C$, the map $s\mapsto \mathcal T^{\mathrm{ren}}(z,s)$ is holomorphic on $\Re s>-1$.
\item The derivative at $s=0$ exists and equals
\[
\frac{\partial}{\partial s}\mathcal T^{\mathrm{ren}}(z,s)\Big|_{s=0}
\ =\ -\sum_{i\ge1}\bigl(\phi_i(z)-\phi_\infty(z)\bigr)\,\log i,
\]
where the series converges absolutely and locally uniformly in $z$.
\item For every integer $n\ge1$ and all $s$ with $\Re s>-1$,
\[
\mathcal T^{\mathrm{ren}}(n,s)=\sum_{d\mid n}d^{-s}=\sigma_{-s}(n).
\]
Consequently,
\[
\frac{\partial}{\partial s}\sigma_{-s}(n)\Big|_{s=0}\ =\ -\sum_{d\mid n}\log d\ =\ -(\tau*\Lambda)(n).
\]
\item On the spectral side,
\[
\sum_{n\ge1}\frac{\sigma_{-s}(n)}{n^{u}}=\zeta(u)\,\zeta(u+s)\qquad(\Re u>1,\ \Re(u+s)>1).
\]
Differentiating at $s=0$ yields
\[
\sum_{n\ge1}\frac{-\sum_{d\mid n}\log d}{n^u}\;=\;\zeta(u)\,\zeta'(u)\qquad(\Re u>1),
\]
and, equivalently,
\[
\sum_{n\ge1}\frac{(\tau*\Lambda)(n)}{n^u}\;=\;-\,\zeta(u)\,\zeta'(u)\qquad(\Re u>1).
\]
Hence, in the Dirichlet–series variable $u$, the differentiation identity is exactly
\[
-\;(\tau*\Lambda)\ \Longleftrightarrow\ \zeta(u)\,\zeta'(u).
\]
\end{enumerate}
\end{proposition}

\begin{lemma}[Divisor–log identity]\label{lem:divlog}
For every $n\ge1$,
\[
\sum_{d\mid n}\log d \;=\; (\tau*\Lambda)(n).
\]
\end{lemma}
\begin{proof}
Two proofs are recorded.

\emph{(i) Dirichlet–series proof.} The Dirichlet series of $n\mapsto \sum_{d\mid n}\log d$ equals
\[
\sum_{n\ge1}\frac{\sum_{d\mid n}\log d}{n^u}
=\sum_{d\ge1}\frac{\log d}{d^{u}}\sum_{m\ge1}\frac{1}{m^{u}}
=\zeta(u)\sum_{d\ge1}\frac{\log d}{d^{u}}
=\zeta(u)\,\bigl(-\zeta'(u)\bigr)\qquad(\Re u>1).
\]
Since $\sum_{n\ge1}(\tau*\Lambda)(n)n^{-u}=\zeta(u)^2\cdot(-\zeta'(u)/\zeta(u))=-\zeta(u)\zeta'(u)$ on $\Re u>1$, both series coincide, hence the coefficients agree.

\emph{(ii) Convolution–algebra proof.} A check on prime powers alone does not suffice here, because neither side is multiplicative: for $n=6$ one has $\sum_{d\mid6}\log d=\log 36=3.5835\ldots$, whereas $\bigl(\sum_{d\mid2}\log d\bigr)\bigl(\sum_{d\mid3}\log d\bigr)=\log2\cdot\log3=0.7615\ldots$. We therefore argue inside the Dirichlet ring. Write $\mathbf 1$ for the constant arithmetic function $1$. By definition $\tau=\mathbf 1*\mathbf 1$, and $\Lambda*\mathbf 1=\log$ is the classical identity $\sum_{d\mid n}\Lambda(d)=\log n$. Associativity and commutativity of Dirichlet convolution give
\[
\tau*\Lambda=(\mathbf 1*\mathbf 1)*\Lambda=\mathbf 1*(\Lambda*\mathbf 1)=\mathbf 1*\log,
\qquad\text{that is,}\qquad
(\tau*\Lambda)(n)=\sum_{d\mid n}\log d .
\]
Pairing each divisor $d$ of $n$ with $n/d$ moreover gives the closed form
\[
2\sum_{d\mid n}\log d=\sum_{d\mid n}\bigl(\log d+\log(n/d)\bigr)=\tau(n)\log n,
\qquad\text{so}\qquad
(\tau*\Lambda)(n)=\tfrac12\,\tau(n)\log n .
\]
\end{proof}

\begin{proof}[Proof of Proposition~\ref{prop:spectral-diff}]
Fix $R>0$ and write $\phi_\infty(z):=(\sin(\pi z)/(\pi z))^2$ with $\phi_\infty(0):=1$. On $|z|\le R$,
\[
\sin(\pi z/i)=\frac{\pi z}{i}+\Big(\frac{\pi z}{i}\Big)^3 \rho_i(z),\qquad |\rho_i(z)|\le \tfrac{1}{6},
\]
whence $\sin(\pi z/i)=\frac{\pi z}{i}\,(1+E_i(z))$ with $|E_i(z)|\le \frac{\pi^2 R^2}{6\,i^{2}}$. Hence
\[
\frac{F(z,i)}{i^2}
=\left(\frac{\sin(\pi z)}{\pi z}\right)^{\!2}\,\frac{1}{(1+E_i(z))^{2}}
=\phi_\infty(z)\,\big(1+O_R(i^{-2})\big),
\]
so $\phi_i(z)-\phi_\infty(z)=O_R(i^{-2})$ uniformly on $|z|\le R$. For $\sigma:=\Re s>-1$,
\[
\sum_{i\ge1}\big|\phi_i(z)-\phi_\infty(z)\big|\,i^{-\sigma}\ \ll_R\ \sum_{i\ge1}i^{-2-\sigma}\ <\ \infty,
\]
locally uniformly in $z$ and on vertical strips $\Re s\ge \sigma_0>-1$. Thus
\[
\mathcal T^{\mathrm{ren}}(z,s):=\sum_{i\ge1}\big(\phi_i(z)-\phi_\infty(z)\big)\,i^{-s}
\]
is entire in $z$ and holomorphic in $s$ for $\Re s>-1$, with termwise differentiation in $s$ justified by $\sum_{i\ge1}(\log i)\,i^{-2-\sigma}$, which converges for $\sigma>-1$. Therefore
\[
\frac{\partial}{\partial s}\mathcal T^{\mathrm{ren}}(z,s)\Big|_{s=0}
=\,-\sum_{i\ge1}\bigl(\phi_i(z)-\phi_\infty(z)\bigr)\,\log i.
\]

At integers $n\ge1$, $\phi_i(n)=\mathbf 1_{i\mid n}$ and $\phi_\infty(n)=0$, so
\[
\mathcal T^{\mathrm{ren}}(n,s)=\sum_{d\mid n}d^{-s}=\sigma_{-s}(n).
\]
Differentiation at $s=0$ gives $\frac{\partial}{\partial s}\sigma_{-s}(n)|_{s=0}=-\sum_{d\mid n}\log d$. Finally,
\[
\sum_{n\ge1}\frac{\sigma_{-s}(n)}{n^u}=\zeta(u)\zeta(u+s)
\quad(\Re u>1,\ \Re(u+s)>1),
\]
so at $s=0$ one obtains $\sum_{n\ge1}\frac{-\sum_{d\mid n}\log d}{n^u}=\zeta(u)\zeta'(u)$.
\end{proof}

\subsubsection{Interpretation and a New Structural Identity}
The transformed spectrum $\zeta(u)\zeta'(u)$ has a known arithmetic counterpart. The Dirichlet series for the divisor-counting function $\tau(n)$ is $\zeta(u)^2$, and for the von Mangoldt function $\Lambda(n)$ it is $-\zeta'(u)/\zeta(u)$. Their product corresponds to the Dirichlet convolution $(\tau * \Lambda)(n)$, and its Dirichlet series is $\zeta(u)^2 \cdot (-\zeta'(u)/\zeta(u)) = -\zeta(u)\zeta'(u)$.

This reveals a structural connection:
\[
\mathcal{D}[\sigma_{-s}(n)] = -\log\left(\prod_{d|n} d\right) \quad \longleftrightarrow \quad \mathcal{D}[\zeta(u)\zeta(u+s)] = \zeta(u)\zeta'(u) = - \sum_{n=1}^\infty \frac{(\tau * \Lambda)(n)}{n^u}.
\]
The differentiation operator provides a constructive analytic bridge from the theory of divisor functions ($\sigma$) to the theory of prime numbers ($\Lambda$), mediated by the divisor-counting function $\tau$. This demonstrates that the relationship between these fundamental arithmetic functions is a direct consequence of the analytic structure of the underlying family of entire functions $\mathcal{T}(z,s)$. The non-obvious arithmetic identity revealed by this principle is visually confirmed in Figure~\ref{fig:operator_effect}.

\begin{figure}[!htbp]
\centering
\includegraphics[width=\textwidth]{plot_operator_effect.pdf}
\caption{A visual representation of the arithmetic identity revealed by the differentiation operator. The plot on the left shows the function $g(n) = -\sum_{d|n} \log d$, derived from differentiating $\sigma_{-s}(n)$. The plot on the right shows the function $h(n) = -(\tau * \Lambda)(n)$, derived from the spectrum $\zeta(u)\zeta'(u)$. Their perfect identity for all $n$ is a visual confirmation of the Principle of Spectral Differentiation.}
\label{fig:operator_effect}
\end{figure}

\FloatBarrier

\paragraph{\textbf{Interpretation.}} The arithmetic identity can be verified algebraically via Dirichlet series. The FD--lift additionally displays how the same relation emerges from an analytic two--variable framework: the differentiation operator in $s$ transports divisor data on the local side to $\zeta(u)\zeta'(u)$ on the spectral side. This view complements the standard convolution calculus without altering its content.

\begin{quote}\small
\textbf{Takeaways (Track 2: dynamic two-variable lift).}
The object $\mathcal T^{(\mathrm{ren})}(z,s)$ is entire in $z$ and holomorphic in $s$ on $\Re s>-1$.
At integers it collapses to $\sum_{d\mid n}d^{-s}$, and the derivative at $s=0$ gives $-\sum_{d\mid n}\log d=-(\tau*\Lambda)(n)$; the von Mangoldt function enters through the spectrum $\zeta(u)\zeta'(u)=\zeta(u)^2\cdot\zeta'(u)/\zeta(u)$, not as the integer value.
Curvature bounds and a Newton template for zero localisation are developed for the indicator $\mathfrak F$ in Application~I (Remark~\ref{rem:newton}).
\end{quote}

% ---------------------------------------------------------------------------
